\section{Introdução}

\paragraph{}
Nesta atividade prática foram exploradas técnicas de enumeração e escalação de
privilégios em ambientes Windows. Diferentemente das atividades anteriores,
onde o levantamento de informações foi realizado predominantemente de forma
manual, neste roteiro foram utilizadas ferramentas especializadas capazes de
automatizar a coleta e análise de informações relevantes para identificação de
vetores de ataque.

\paragraph{}
Durante a execução das atividades foram analisadas configurações de serviços,
permissões inadequadas, problemas relacionados a caminhos de executáveis sem
aspas, carregamento inseguro de bibliotecas dinâmicas e exposição de arquivos
sensíveis do sistema operacional.

\paragraph{}
Além da utilização da ferramenta \texttt{WinPEAS} para enumeração do ambiente,
foram empregadas ferramentas auxiliares para validação de permissões,
monitoramento de processos, criação de binários compatíveis com serviços do
Windows e recuperação de credenciais armazenadas localmente.

\paragraph{}
O objetivo principal da prática foi compreender como falhas de configuração e
controles inadequados de segurança podem ser explorados para obtenção de
privilégios elevados em sistemas Windows, reforçando a importância da correta
administração e proteção desses ambientes.
